\documentclass{article}

\usepackage[utf8]{inputenc}
\usepackage[T1]{fontenc}
\usepackage[spanish]{babel}
\usepackage{amsmath}
\usepackage{amssymb}
\usepackage{array}
\usepackage{geometry}
\usepackage{enumitem}

\geometry{margin=2.5cm}

\begin{document}

\section{Análisis léxico con Flex}

El análisis léxico puede modelarse conceptualmente como una función que
recorre el código fuente, identifica lexemas y los transforma en tokens.
Flex permite automatizar la implementación de esta función a partir de una
especificación de patrones léxicos.

Sea $C$ el código fuente y sea $T$ el conjunto de tokens del lenguaje.
Podemos representar el analizador léxico como:

\[
L : C \rightarrow T^*
\]

donde $T^*$ representa una secuencia de tokens.

De forma recursiva, el lexer puede interpretarse como una función que
consume progresivamente la entrada. Si $c$ representa la entrada restante,
el analizador busca el siguiente lexema reconocido por alguno de los
patrones definidos y ejecuta la acción asociada.

\[
L(c) =
\begin{cases}
[] & \text{si } c = \epsilon \\[4pt]
t :: L(c') & \text{si el lexema } l \text{ coincide con un patrón y}\\
& \text{la acción asociada produce el token } t \\[4pt]
l :: L(c') & \text{si el lexema } l \text{ no coincide con ningún patrón}
\end{cases}
\]

donde $c'$ representa la entrada restante después de consumir el lexema
$l$.

\subsection{Patrones y acciones}

Una especificación de Flex está formada por reglas que relacionan un
\textit{patrón} con una \textit{acción}. El comienzo de esta sección se
indica mediante \texttt{\%\%}.

Por ejemplo:

\begin{verbatim}
%%
"username"    printf("%s", getlogin());
\end{verbatim}

En esta regla, \texttt{"username"} es el patrón y
\texttt{printf("\%s", getlogin())} es la acción.

Cuando Flex encuentra la cadena \texttt{username} en la entrada, ejecuta
la acción asociada y reemplaza dicha cadena por el nombre de usuario
obtenido mediante \texttt{getlogin()}.

Por defecto, cuando Flex encuentra texto que no coincide con ninguno de
los patrones definidos, dicho texto se copia directamente a la salida.
Por lo tanto, en el ejemplo anterior, el analizador copia el archivo de
entrada a la salida, excepto por las apariciones de \texttt{"username"},
que son reemplazadas por el nombre de usuario correspondiente.

Este comportamiento permite entender una regla de Flex como una
transformación:

\[
\text{patrón léxico}
\longrightarrow
\text{acción}
\]

donde la acción determina qué ocurre cuando el patrón es reconocido.

\subsection{Expresiones regulares y patrones de Flex}

Los patrones de Flex se construyen principalmente mediante expresiones
regulares. Estas permiten describir conjuntos de cadenas que pueden
aparecer en el código fuente.

\subsubsection{Caracteres literales}

Un carácter escrito directamente en el patrón representa ese carácter.

\begin{verbatim}
+
-
*
=
;
\end{verbatim}

Por ejemplo:

\begin{verbatim}
"+"    { return SUMA; }
"="    { return ASIGNACION; }
";"    { return PUNTO_COMA; }
\end{verbatim}

Las comillas permiten representar explícitamente una cadena literal.

\subsubsection{Cadenas literales}

Una secuencia de caracteres entre comillas representa exactamente esa
secuencia.

\begin{verbatim}
"int"
"while"
"return"
"=="
\end{verbatim}

Por ejemplo:

\begin{verbatim}
"while"    { return WHILE; }
"return"   { return RETURN; }
\end{verbatim}

\subsubsection{Clase de caracteres}

Los corchetes permiten especificar un conjunto de caracteres. El patrón
coincide con exactamente un carácter perteneciente a ese conjunto.

\begin{verbatim}
[abc]
\end{verbatim}

coincide con:

\begin{verbatim}
a
b
c
\end{verbatim}

También pueden utilizarse rangos:

\begin{verbatim}
[a-z]
[A-Z]
[0-9]
\end{verbatim}

Por ejemplo:

\begin{verbatim}
[0-9]    { return DIGITO; }
\end{verbatim}

reconoce cualquier dígito entre $0$ y $9$.

Los rangos pueden combinarse:

\begin{verbatim}
[a-zA-Z]
[a-zA-Z0-9]
\end{verbatim}

\subsubsection{Negación de una clase}

El carácter \texttt{\^{}} como primer carácter dentro de una clase niega
la clase.

\begin{verbatim}
[^0-9]
\end{verbatim}

coincide con cualquier carácter que no sea un dígito.

Por ejemplo:

\begin{verbatim}
[^a-zA-Z0-9]
\end{verbatim}

coincide con cualquier carácter que no sea una letra ni un número.

\subsubsection{Concatenación}

Dos patrones escritos uno después del otro representan una concatenación.
Ambos deben aparecer consecutivamente en la entrada.

Por ejemplo:

\begin{verbatim}
ab
\end{verbatim}

coincide con:

\begin{verbatim}
ab
\end{verbatim}

pero no con:

\begin{verbatim}
a
b
ac
\end{verbatim}

Un identificador puede expresarse mediante:

\begin{verbatim}
[a-zA-Z_][a-zA-Z0-9_]*
\end{verbatim}

donde primero debe aparecer una letra o guion bajo y posteriormente
pueden aparecer letras, números o guiones bajos.

\subsubsection{Alternancia}

El operador \texttt{|} representa una alternativa. El patrón coincide con
una opción u otra.

\begin{verbatim}
if|while|for
\end{verbatim}

coincide con cualquiera de:

\begin{verbatim}
if
while
for
\end{verbatim}

También puede utilizarse con grupos:

\begin{verbatim}
(if|while|for)
\end{verbatim}

\subsubsection{Cero o más repeticiones}

El operador \texttt{*} significa que el elemento anterior puede aparecer
cero o más veces.

\begin{verbatim}
[0-9]*
\end{verbatim}

coincide con:

\begin{verbatim}
""
"1"
"123"
"98765"
\end{verbatim}

Por ejemplo:

\begin{verbatim}
[a-zA-Z_][a-zA-Z0-9_]*
\end{verbatim}

permite que después del primer carácter de un identificador aparezca
cualquier cantidad de letras, números o guiones bajos.

\subsubsection{Una o más repeticiones}

El operador \texttt{+} significa que el elemento anterior debe aparecer
una o más veces.

\begin{verbatim}
[0-9]+
\end{verbatim}

coincide con:

\begin{verbatim}
1
123
98765
\end{verbatim}

pero no coincide con una cadena vacía.

Por ejemplo, esta es una forma habitual de reconocer números enteros:

\begin{verbatim}
[0-9]+    { return NUMERO; }
\end{verbatim}

\subsubsection{Cero o una aparición}

El operador \texttt{?} indica que el elemento anterior puede aparecer
cero o una vez.

\begin{verbatim}
-?[0-9]+
\end{verbatim}

permite reconocer:

\begin{verbatim}
123
-123
45
-45
\end{verbatim}

Por lo tanto, puede utilizarse para representar un signo negativo
opcional.

\subsubsection{Repeticiones exactas y rangos}

Las expresiones entre llaves permiten especificar cantidades de
repeticiones.

\begin{verbatim}
[0-9]{4}
\end{verbatim}

indica exactamente cuatro dígitos.

Por ejemplo:

\begin{verbatim}
2026
1234
9876
\end{verbatim}

También pueden especificarse rangos:

\begin{verbatim}
[0-9]{2,4}
\end{verbatim}

que permite entre dos y cuatro dígitos.

\begin{verbatim}
[0-9]{2,}
\end{verbatim}

permite dos o más dígitos.

\subsubsection{Cualquier carácter}

El punto \texttt{.} representa cualquier carácter excepto, en general,
un salto de línea.

Por ejemplo:

\begin{verbatim}
.    { return CARACTER; }
\end{verbatim}

puede utilizarse como regla general para detectar caracteres que no fueron
reconocidos por reglas anteriores.

Esta característica resulta especialmente útil para detectar errores
léxicos:

\begin{verbatim}
. {
    printf("Error léxico: %s\n", yytext);
}
\end{verbatim}

\subsubsection{Paréntesis y agrupación}

Los paréntesis permiten agrupar patrones para aplicar operadores sobre
varios elementos.

Por ejemplo:

\begin{verbatim}
(ab)+
\end{verbatim}

coincide con:

\begin{verbatim}
ab
abab
ababab
\end{verbatim}

También pueden combinarse con alternancia:

\begin{verbatim}
(if|while|for)
\end{verbatim}

\subsubsection{Caracteres especiales}

Algunos caracteres poseen un significado especial dentro de una
expresión regular. Entre ellos se encuentran:

\begin{verbatim}
.  *  +  ?  |  (  )  [  ]  {  }  ^  $
\end{verbatim}

Cuando se necesita reconocer literalmente uno de estos caracteres,
puede utilizarse una barra inversa para escaparlo.

Por ejemplo:

\begin{verbatim}
\+
\*
\(
\)
\[
\]
\{
\}
\.
\end{verbatim}

También es habitual utilizar comillas para representar operadores como
literales:

\begin{verbatim}
"+"     { return SUMA; }
"*"     { return MULTIPLICACION; }
"("     { return PARENTESIS_IZQ; }
")"     { return PARENTESIS_DER; }
\end{verbatim}

\subsubsection{Caracteres especiales de escape}

Flex permite utilizar secuencias de escape para representar caracteres
especiales.

Algunos de los más utilizados son:

\begin{center}
\begin{tabular}{c|l}
\textbf{Secuencia} & \textbf{Significado} \\
\hline
\texttt{\textbackslash n} & salto de línea \\
\texttt{\textbackslash t} & tabulación \\
\texttt{\textbackslash r} & retorno de carro \\
\texttt{\textbackslash f} & salto de página \\
\texttt{\textbackslash v} & tabulación vertical \\
\texttt{\textbackslash \textbackslash} & barra inversa \\
\end{tabular}
\end{center}

Por ejemplo:

\begin{verbatim}
[ \t\n]+    { /* ignorar espacios */ }
\end{verbatim}

\subsubsection{Inicio y final de línea}

El carácter \texttt{\^} fuera de una clase de caracteres puede utilizarse
para indicar el comienzo de una línea.

Por ejemplo:

\begin{verbatim}
^int
\end{verbatim}

coincide con \texttt{int} cuando aparece al comienzo de una línea.

El carácter \texttt{\$} puede utilizarse para indicar el final de una
línea.

Por ejemplo:

\begin{verbatim}
int$
\end{verbatim}

coincide con \texttt{int} cuando aparece al final de una línea.

\subsubsection{Contexto posterior}

Flex permite utilizar el operador \texttt{/} para especificar contexto
posterior.

Por ejemplo:

\begin{verbatim}
[0-9]+/[a-zA-Z]
\end{verbatim}

indica que una secuencia de números debe estar seguida por una letra para
que el patrón sea considerado un emparejamiento válido.

El texto utilizado como contexto posterior participa en la determinación
del emparejamiento, aunque no necesariamente forma parte del texto
devuelto como token.

\subsubsection{Macros de patrones}

Flex permite definir nombres para patrones en la sección de
definiciones. Esto evita repetir expresiones regulares complejas.

Por ejemplo:

\begin{verbatim}
DIGITO      [0-9]
LETRA       [a-zA-Z]
ID          {LETRA}({LETRA}|{DIGITO})*
\end{verbatim}

Luego pueden utilizarse estos nombres en las reglas:

\begin{verbatim}
%%
{ID}       { return IDENTIFICADOR; }
{DIGITO}+  { return NUMERO; }
\end{verbatim}

Esto permite construir patrones complejos a partir de patrones más
simples y hace que la especificación sea más legible.

\subsection{Resumen de patrones}

Los principales operadores de expresiones regulares utilizados en Flex
pueden resumirse de la siguiente manera:

\begin{center}
\begin{tabular}{c|l}
\textbf{Patrón} & \textbf{Significado} \\
\hline
\texttt{a} & carácter literal \\
\texttt{"abc"} & cadena literal \\
\texttt{[abc]} & uno de los caracteres indicados \\
\texttt{[a-z]} & un carácter dentro de un rango \\
\texttt{[\^{}a-z]} & cualquier carácter fuera del rango \\
\texttt{.} & cualquier carácter excepto salto de línea \\
\texttt{ab} & concatenación de patrones \\
\texttt{a|b} & alternativa entre patrones \\
\texttt{a*} & cero o más repeticiones \\
\texttt{a+} & una o más repeticiones \\
\texttt{a?} & cero o una aparición \\
\texttt{a\{n\}} & exactamente $n$ repeticiones \\
\texttt{a\{n,m\}} & entre $n$ y $m$ repeticiones \\
\texttt{(ab)} & agrupación \\
\texttt{\^{}a} & comienzo de línea \\
\texttt{a\$} & final de línea \\
\texttt{a/b} & contexto posterior \\
\texttt{\textbackslash c} & carácter especial escapado \\
\end{tabular}
\end{center}

Estos patrones constituyen las herramientas fundamentales para describir
los lexemas de un lenguaje de programación.

\subsection{Proceso de emparejamiento}

Cuando el escáner generado por Flex se encuentra en una determinada
posición de la entrada, busca los patrones que pueden comenzar en dicha
posición. Cada patrón intenta consumir la mayor cantidad de caracteres
posible.

Por ejemplo, si se definen las siguientes reglas:

\begin{verbatim}
[0-9]     { return DIGITO; }
[0-9]+    { return NUMERO; }
\end{verbatim}

y la entrada es:

\begin{verbatim}
1234 + 1
\end{verbatim}

al comenzar en el primer carácter, ambas reglas pueden coincidir:

\[
\begin{array}{c|c}
\text{Patrón} & \text{Emparejamiento} \\
\hline
\texttt{[0-9]} & \texttt{1} \\
\texttt{[0-9]+} & \texttt{1234}
\end{array}
\]

Flex selecciona el emparejamiento que consume más caracteres. Por lo
tanto, selecciona:

\[
\texttt{1234} \longrightarrow \text{NUMERO}
\]

y no:

\[
\texttt{1} \longrightarrow \text{DIGITO}
\]

La posición de análisis avanza solamente después de haber determinado
el emparejamiento ganador. En este caso, después de reconocer
\texttt{1234}, la entrada restante es:

\begin{verbatim}
+ 1
\end{verbatim}

A continuación, Flex repite el mismo proceso desde la nueva posición.

De esta manera, el proceso puede resumirse en los siguientes pasos:

\begin{enumerate}
\item Se fija la posición actual de la entrada.
\item Se buscan todos los patrones que pueden coincidir desde esa
posición.
\item Cada patrón intenta extender el emparejamiento tanto como sea
posible.
\item Se conservan los emparejamientos válidos encontrados.
\item Se selecciona el emparejamiento que consume la mayor cantidad
de caracteres.
\item Se ejecuta la acción asociada al patrón seleccionado.
\item La posición de la entrada avanza hasta el final del lexema
reconocido.
\item El proceso se repite hasta alcanzar el final de la entrada.
\end{enumerate}

Este comportamiento se conoce como el principio de
\textit{emparejamiento más largo} o \textit{maximal munch}.

Si dos o más patrones producen emparejamientos de la misma longitud,
Flex selecciona el patrón que aparece primero en el archivo de
especificación.

\subsection{Flex como productor de tokens}

Cuando Flex se utiliza como parte de un compilador, las acciones de las
reglas pueden utilizarse para producir tokens que posteriormente serán
consumidos por el parser generado por Bison.

Por ejemplo:

\begin{verbatim}
%%
"int"                         { return INT; }
[0-9]+                        { return NUMERO; }
[a-zA-Z_][a-zA-Z0-9_]*        { return IDENTIFICADOR; }
"+"                           { return SUMA; }
"="                           { return ASIGNACION; }
";"                           { return PUNTO_COMA; }
[ \t\n]+                      { /* ignorar */ }
\end{verbatim}

Para una entrada como:

\begin{verbatim}
int x = 10 + 20;
\end{verbatim}

el análisis léxico produce conceptualmente:

\[
[
\text{INT},
\text{IDENTIFICADOR},
\text{ASIGNACION},
\text{NUMERO},
\text{SUMA},
\text{NUMERO},
\text{PUNTO\_COMA}
]
\]

Los espacios, tabulaciones y saltos de línea pueden ser reconocidos por
una regla específica y descartados mediante una acción vacía.

\subsection{Lexema reconocido y variables de Flex}

Una vez determinado el emparejamiento, Flex pone a disposición de la
acción correspondiente el texto reconocido mediante la variable
\texttt{yytext} y su longitud mediante la variable global entera
\texttt{yyleng}.

Por ejemplo:

\begin{verbatim}
[0-9]+ {
    printf("Lexema: %s\n", yytext);
    printf("Longitud: %d\n", yyleng);
}
\end{verbatim}

Si la entrada contiene:

\begin{verbatim}
1234
\end{verbatim}

el valor de \texttt{yytext} será:

\begin{verbatim}
1234
\end{verbatim}

mientras que \texttt{yyleng} tendrá el valor:

\[
4
\]

Esto permite que la acción asociada a una regla trabaje con el lexema
completo que fue reconocido y no solamente con el primer carácter que
coincidió con el patrón.

\subsection{Camino no feliz}

Un caso no feliz ocurre cuando el código fuente contiene un carácter o
secuencia que no coincide con ninguno de los patrones definidos.

Por ejemplo:

\begin{verbatim}
int x = 10 @ 20;
\end{verbatim}

Si el lenguaje no define el carácter \texttt{@}, podemos utilizar una
regla general para detectarlo:

\begin{verbatim}
. {
    printf("Error léxico: %s\n", yytext);
}
\end{verbatim}

El procesamiento será conceptualmente:

\[
\begin{array}{c|c}
\text{Lexema} & \text{Resultado} \\
\hline
\texttt{int} & \text{INT} \\
\texttt{x} & \text{IDENTIFICADOR} \\
\texttt{=} & \text{ASIGNACION} \\
\texttt{10} & \text{NUMERO} \\
\texttt{@} & \text{ERROR LÉXICO} \\
\texttt{20} & \text{NUMERO} \\
\texttt{;} & \text{PUNTO\_COMA}
\end{array}
\]

Por lo tanto, Flex recorre la entrada buscando continuamente el
emparejamiento más largo disponible en cada posición. Una vez reconocido
el lexema, ejecuta la acción correspondiente y continúa el análisis desde
el carácter siguiente.

Finalmente, el flujo completo puede resumirse como:

\[
\text{Código fuente}
\xrightarrow{\text{Flex}}
\text{Secuencia de tokens}
\xrightarrow{\text{Bison}}
\text{Análisis sintáctico}
\]

\end{document}